\documentclass[conference]{IEEEtran}
\IEEEoverridecommandlockouts

\usepackage{cite}
\usepackage{url}
\usepackage{amsmath}
\usepackage{booktabs}
\usepackage{array}

\title{AegisMind: Defense-in-Depth Security Framework for Agentic AI Workflows (Current Version)}

\author{
\IEEEauthorblockN{A. Doni Adithya, E. Hari Teja, Hari Prasad L., L. Sahith Akash Manikanta}
\IEEEauthorblockA{
Registration Numbers: 23bds007, 23bds019, 23bds022, 23bds031 \\
Course: Data Security and Privacy (DS-208) \\
Professor: Dr Girish Revadigar
}
}

\begin{document}
\maketitle

\begin{abstract}
Agentic AI systems improve automation by combining language reasoning with tool execution, but this increases the risk of prompt injection, role abuse, unsafe tool usage, data leakage, and operational overload. This paper presents the current implemented version of \textit{AegisMind}, a FastAPI and React based defense-in-depth framework for secure AI-agent execution. The system combines multi-model prompt risk analysis (regex + machine learning + LLM), role-aware policy enforcement, per-tool risk profiles, human approval workflow, secure gateway mediation, DLP output controls, behavioral remediation, encrypted auditability, and operational governance capabilities. The latest version adds dashboard user authentication, policy versioning and change audit, threat-intelligence feed sync, backup/restore operations, scheduled key rotation and log archival controls, realtime event streaming, and scorecard analytics. Experimental checks show stable backend and frontend quality gates in the current build.
\end{abstract}

\begin{IEEEkeywords}
Agentic AI security, prompt firewall, policy governance, DLP, human approval, threat intelligence, observability.
\end{IEEEkeywords}

\section{Introduction}
As LLM systems evolve from chat interfaces to action-capable agents, untrusted text becomes a control interface over tools, APIs, and data. This introduces direct security risks: malicious prompt steering, privilege escalation through natural language, and accidental sensitive-data disclosure via tool output.

AegisMind is designed as a runtime security control plane for such agents. Instead of relying on a single detector, it applies layered controls from request ingress to final response.

\section{Problem Definition and Threat Model}
\subsection{Problem Statement}
Given a prompt and optional tool request, determine whether execution should be allowed, modified, approval-gated, or blocked while preserving confidentiality, integrity, and availability.

\subsection{Threat Assumptions}
Adversaries may attempt:
\begin{itemize}
    \item prompt injection and jailbreak instructions,
    \item role or policy manipulation,
    \item encoded malicious intent (for example Base64),
    \item unauthorized tool execution,
    \item secret/PII extraction through output channels,
    \item request floods causing performance degradation.
\end{itemize}

\section{Current Architecture}
\subsection{Execution Pipeline}
\begin{enumerate}
    \item Request enters \texttt{/agent/execute}.
    \item Identity is verified (agent key headers or dashboard bearer token).
    \item Rate limiting and queue/backpressure constraints are applied.
    \item Prompt firewall returns risk score and decision.
    \item Tool access is evaluated via role policy and per-tool profile.
    \item Human approval is required for elevated-risk operations.
    \item Secure gateway executes allow-listed tools only.
    \item DLP scans and masks/blocks sensitive output.
    \item Audit and security events are persisted (encrypted text fields).
    \item Realtime event updates are broadcast via WebSocket.
\end{enumerate}

\subsection{Implemented Core Modules}
\begin{itemize}
    \item Prompt Firewall with multi-model voting metadata.
    \item Policy Engine with role rules, global denies, unsafe chain checks.
    \item Tool Risk Profiles and Human Approval workflow.
    \item DLP for output sanitization and hard blocking.
    \item Behavioral self-healing states.
    \item Threat-intelligence import/sync integration.
    \item Operations services (archive, backup/restore, key rotation, SLO endpoint).
\end{itemize}

\section{Implementation Details (Current Version)}
\subsection{Authentication and Governance}
The current version supports dashboard login endpoints (\texttt{/auth/login}, \texttt{/auth/logout}, \texttt{/auth/me}) and admin user management (\texttt{/auth/users}) with role and team metadata.

Policy governance includes:
\begin{itemize}
    \item policy export/apply and publish endpoints,
    \item policy version persistence,
    \item ``who changed what'' policy change audit trail.
\end{itemize}

\subsection{Prompt Risk and Explainability}
The firewall exposes explainability fields including matched rules, threat signals, multi-model component scores, and calibration metadata. Risk zoning follows:
\begin{itemize}
    \item Safe Zone: 0--60\%
    \item Block Zone: 60--100\%
\end{itemize}

\subsection{Operations and Reliability}
The system includes:
\begin{itemize}
    \item queue and semaphore-based concurrency controls,
    \item timeout budgets and backpressure responses,
    \item webhook/Slack/email alert fan-out support,
    \item archive listing and retention endpoints,
    \item backup create/list/restore endpoints,
    \item calibration feedback and bias summary APIs,
    \item SLO status endpoint and optional OpenTelemetry bootstrap.
\end{itemize}

\subsection{Threat Intelligence and Realtime}
Threat patterns can be imported manually or synchronized from a remote feed. A scheduled maintenance loop can periodically refresh feeds and apply maintenance tasks. Realtime stream updates are available via \texttt{/ws/security-stream}.

\section{Frontend System}
The React dashboard currently includes:
\begin{itemize}
    \item \textbf{Dashboard}: simulator, explainability, risk zone summary.
    \item \textbf{Agent}: reasoning trace and execution inspection.
    \item \textbf{Logs}: searchable audit stream and exports.
    \item \textbf{Settings}: runtime connection controls.
    \item \textbf{Ops}: policy studio, approvals, threat feed controls, tool profiles, replay, realtime stream, scorecard, calibration, backups, archives, key rotation, and user management.
\end{itemize}

\section{Evaluation Snapshot}
Table~\ref{tab:validation} summarizes the latest validation results for the current codebase.

\begin{table}[h]
\centering
\caption{Current Validation Results}
\label{tab:validation}
\begin{tabular}{>{\raggedright\arraybackslash}p{3.0cm}p{3.7cm}}
\toprule
\textbf{Check} & \textbf{Result} \\
\midrule
Backend tests (\texttt{pytest}) & 18 passed \\
Frontend lint (\texttt{npm run lint}) & passed \\
Frontend build (\texttt{npm run build}) & passed \\
Frontend e2e (\texttt{playwright}) & 3 passed \\
\bottomrule
\end{tabular}
\end{table}

\section{Discussion}
The current version demonstrates that agent security improves significantly when detection, authorization, and operations are integrated. Multi-model scoring provides robust threat sensitivity, while policy, tool profiles, and approval gates reduce blast radius for risky actions. Operational APIs make the system practical for continuous monitoring and governance, not just one-time filtering.

\section{Limitations}
\begin{itemize}
    \item Optional OpenTelemetry depends on deployment-specific package availability.
    \item Team metadata exists, but full enterprise IAM federation is not yet integrated.
    \item Threat feed quality depends on external source quality and curation.
    \item Calibration currently relies on operator feedback, which can be sparse.
\end{itemize}

\section{Future Work}
\begin{itemize}
    \item Integrate enterprise SSO (OIDC/SAML) and fine-grained team scopes.
    \item Add policy simulation tests before publish (pre-merge guardrails).
    \item Extend SLO telemetry into Grafana/Prometheus dashboards.
    \item Expand adversarial benchmark corpus under high-concurrency traffic.
\end{itemize}

\section{Conclusion}
The current AegisMind release provides a complete, test-validated defense-in-depth platform for agentic AI security. It now combines runtime threat controls, policy governance, approval workflows, threat-intelligence ingestion, operational resiliency features, and actionable observability in one integrated architecture. This aligns with DS-208 goals by converting privacy and security principles into an operational and auditable system.

\section*{Acknowledgment}
This project was completed for \textbf{Data Security and Privacy (DS-208)} under the guidance of \textbf{Dr Girish Revadigar}.

\begin{thebibliography}{99}
\bibitem{fastapi} FastAPI Documentation. [Online]. Available: \url{https://fastapi.tiangolo.com/}
\bibitem{sqlalchemy} SQLAlchemy Documentation. [Online]. Available: \url{https://docs.sqlalchemy.org/}
\bibitem{redis} Redis Documentation. [Online]. Available: \url{https://redis.io/docs/}
\bibitem{docker} Docker Documentation. [Online]. Available: \url{https://docs.docker.com/}
\bibitem{sklearn} Scikit-learn Documentation. [Online]. Available: \url{https://scikit-learn.org/stable/}
\bibitem{pydantic} Pydantic Documentation. [Online]. Available: \url{https://docs.pydantic.dev/}
\bibitem{cryptography} Cryptography Documentation. [Online]. Available: \url{https://cryptography.io/}
\bibitem{owasp} OWASP Top 10 for LLM Applications. [Online]. Available: \url{https://owasp.org/www-project-top-10-for-large-language-model-applications/}
\end{thebibliography}

\end{document}
